\hypertarget{the-python-data-model-comprehensive-dunder-methods}{%
\chapter{The Python Data Model \& Comprehensive Dunder
Methods}\label{the-python-data-model-comprehensive-dunder-methods}}

The ``Data Model'' is the formal description of Python's objects and
their interactions. While most developers know
\passthrough{\lstinline!\_\_init\_\_!}, the ``Godhood'' level of
understanding requires knowing how these high-level methods map directly
to functional pointers in the CPython C source code.

\hypertarget{the-philosophy-protocols-over-types}{%
\subsection{32.1 The Philosophy: Protocols over
Types}\label{the-philosophy-protocols-over-types}}

Python uses ``duck typing,'' but it is more accurately described as a
\textbf{Protocol-based language}. If an object implements the methods
required by a protocol, it \emph{is} that thing. These protocols are
implemented using \textbf{Special Methods} (Dunder methods).

\hypertarget{object-lifecycle-and-representation}{%
\subsection{32.2 Object Lifecycle and
Representation}\label{object-lifecycle-and-representation}}

\hypertarget{creation-and-initialization}{%
\subsubsection{1. Creation and
Initialization}\label{creation-and-initialization}}

\begin{itemize}
\tightlist
\item
  \passthrough{\lstinline!\_\_new\_\_(cls, ...)!}: The actual
  constructor. It returns a new instance of
  \passthrough{\lstinline!cls!}. It maps to the
  \passthrough{\lstinline!tp\_new!} slot in C.
\item
  \passthrough{\lstinline!\_\_init\_\_(self, ...)!}: The initializer. It
  configures the instance created by
  \passthrough{\lstinline!\_\_new\_\_!}. It maps to
  \passthrough{\lstinline!tp\_init!}.
\end{itemize}

\hypertarget{string-representations}{%
\subsubsection{2. String Representations}\label{string-representations}}

\begin{itemize}
\tightlist
\item
  \passthrough{\lstinline!\_\_repr\_\_(self)!}: The ``official'' string
  representation, ideally usable to recreate the object. Used by the
  debugger and REPL. Maps to \passthrough{\lstinline!tp\_repr!}.
\item
  \passthrough{\lstinline!\_\_str\_\_(self)!}: The ``informal'' or
  user-friendly string representation. Maps to
  \passthrough{\lstinline!tp\_str!}.
\end{itemize}

\hypertarget{the-mapping-to-c-slots}{%
\subsection{32.3 The Mapping to C Slots}\label{the-mapping-to-c-slots}}

Every Python class is an instance of
\passthrough{\lstinline!PyTypeObject!}. This C struct contains a vast
array of ``slots''function pointers that the VM calls when executing
operations.

\begin{longtable}[]{@{}lll@{}}
\toprule
Python Method & C Slot & Description \\
\midrule
\endhead
\passthrough{\lstinline!\_\_call\_\_!} &
\passthrough{\lstinline!tp\_call!} & Called when object is invoked like
a function. \\
\passthrough{\lstinline!\_\_iter\_\_!} &
\passthrough{\lstinline!tp\_iter!} & Returns an iterator object. \\
\passthrough{\lstinline!\_\_next\_\_!} &
\passthrough{\lstinline!tp\_iternext!} & Returns the next item from an
iterator. \\
\passthrough{\lstinline!\_\_getattr\_\_!} &
\passthrough{\lstinline!(dynamic)!} & Called if attribute lookup
fails. \\
\passthrough{\lstinline!\_\_getattribute\_\_!} &
\passthrough{\lstinline!tp\_getattro!} & Called for EVERY attribute
lookup. \\
\bottomrule
\end{longtable}

\hypertarget{numeric-and-container-protocols}{%
\subsection{32.4 Numeric and Container
Protocols}\label{numeric-and-container-protocols}}

To save space and optimize lookup, CPython groups related methods into
sub-structs:

\hypertarget{tp_as_number}{%
\subsubsection{\texorpdfstring{1.
\texttt{tp\_as\_number}}{1. tp\_as\_number}}\label{tp_as_number}}

Methods like \passthrough{\lstinline!\_\_add\_\_!},
\passthrough{\lstinline!\_\_sub\_\_!}, and
\passthrough{\lstinline!\_\_mul\_\_!} are stored here. When you write
\passthrough{\lstinline!a + b!}, the VM looks at
\passthrough{\lstinline!a->ob\_type->tp\_as\_number->nb\_add!}.

\hypertarget{tp_as_sequence-and-tp_as_mapping}{%
\subsubsection{\texorpdfstring{2. \texttt{tp\_as\_sequence} and
\texttt{tp\_as\_mapping}}{2. tp\_as\_sequence and tp\_as\_mapping}}\label{tp_as_sequence-and-tp_as_mapping}}

\begin{itemize}
\tightlist
\item
  \textbf{Sequence}: \passthrough{\lstinline!\_\_len\_\_!}
  (\passthrough{\lstinline!sq\_length!}),
  \passthrough{\lstinline!\_\_getitem\_\_!}
  (\passthrough{\lstinline!sq\_item!}).
\item
  \textbf{Mapping}: \passthrough{\lstinline!\_\_getitem\_\_!}
  (\passthrough{\lstinline!mp\_subscript!}),
  \passthrough{\lstinline!\_\_setitem\_\_!}
  (\passthrough{\lstinline!mp\_ass\_subscript!}).
\end{itemize}

Note that \passthrough{\lstinline!\_\_getitem\_\_!} is overloaded; if
the object is a sequence, it expects an integer; if it's a mapping, it
expects a hashable key.

\hypertarget{comprehensive-comparison-rich-comparisons}{%
\subsection{32.5 Comprehensive Comparison: Rich
Comparisons}\label{comprehensive-comparison-rich-comparisons}}

Python 3 unified comparisons into \textbf{Rich Comparisons}
(\passthrough{\lstinline!tp\_richcompare!}). *
\passthrough{\lstinline!\_\_lt\_\_!},
\passthrough{\lstinline!\_\_le\_\_!},
\passthrough{\lstinline!\_\_eq\_\_!},
\passthrough{\lstinline!\_\_ne\_\_!},
\passthrough{\lstinline!\_\_gt\_\_!},
\passthrough{\lstinline!\_\_ge\_\_!}.

These all map to a single C function that receives an
\passthrough{\lstinline!op!} argument (e.g.,
\passthrough{\lstinline!Py\_EQ!}, \passthrough{\lstinline!Py\_LT!}). If
you implement only \passthrough{\lstinline!\_\_eq\_\_!}, Python does not
automatically derive \passthrough{\lstinline!\_\_ne\_\_!} or others,
unlike some older versions.

\hypertarget{the-__slots__-optimization-recap}{%
\subsection{\texorpdfstring{32.6 The \texttt{\_\_slots\_\_} Optimization
(Recap)}{32.6 The \_\_slots\_\_ Optimization (Recap)}}\label{the-__slots__-optimization-recap}}

As covered in Chapter 26, \passthrough{\lstinline!\_\_slots\_\_!}
prevents the creation of \passthrough{\lstinline!\_\_dict\_\_!}.
Internally, this changes the \passthrough{\lstinline!PyTypeObject!}
flags and allocates space for the attributes directly in the object
struct, mapping them via \textbf{Member Descriptors}.

\begin{center}\rule{0.5\linewidth}{0.5pt}\end{center}

\begin{center}\rule{0.5\linewidth}{0.5pt}\end{center}
